\section{Conclusion}\label{sec:conc}
This paper presented an evidence-level routing system for structured
aerial VQA that jointly selects the transmitted representation and
the receiver's answering computation. A question-type rule exploits
differences in evidence requirements, while per-sample correctness
prediction and an energy price adjust the accuracy--energy trade-off.
On VisDrone, the type rule improves accuracy over fixed-image
transmission and reduces accounted system energy by about $55\%$
across three channel models. On DroneVehicle, the validation-selected
MLP improves accuracy by $1.67$ percentage points and reduces
accounted system energy by $60.1\%$ relative to the type rule.
The benefit of learned routing depends on the dataset and receiver;
nonlinear prediction is not uniformly better than a linear baseline.
These findings support matching both communication content and
receiver computation to the question, with energy savings arising
mainly from avoiding VLM inference when detection evidence is sufficient.
The savings apply to the evaluated UAV-and-edge energy model and
do not imply lower UAV energy, since onboard detection can outweigh
radio savings. Future work will evaluate end-to-end energy and
latency for online routing on embedded platforms.
